\section{Conclusions}
\label{sec:conclusions}

This work presented a CFD-ML methodology to build a turbulence closure aimed at improving the fidelity of
RANS simulations of HAWT wakes while keeping the
computational cost compatible with wind farm applications. The ML
module is used to define the turbulence model itself, by
learning the eddy-viscosity correction required to reproduce
high-fidelity LES statistics.

A LES of the NREL-5MW turbine operating in a turbulent atmospheric
boundary layer was used to generate the reference data. An
LES-equivalent eddy viscosity $\nu_{t,eq}$ was extracted by projecting
the deviatoric LES stress onto the Boussinesq form, and a Multi-Layer
Perceptron was trained to predict the local correction
$\Delta\nu_t = \nu_{t,eq} - \nu_{t,\mathrm{RANS}}$ with respect to the
RANS baseline. A K-means clustering algorithm was employed to segment
the domain into physically distinct flow regions, moving the prediction
from a cell-wise to a region-aware basis, and the resulting closure was
embedded into a modified OpenFOAM PIMPLE solver to be used in a fully
predictive setting. The main findings can be summarised as follows.

\begin{itemize}
    \item The clustering identified seven physically meaningful flow
    regions, separating the undisturbed flow from the
    inlet/ABL-affected, near-wake, mixing/breakdown and far-wake regions,
    consistently with the underlying flow physics.

    \item The trained network reproduced the LES-equivalent eddy
    viscosity with a validation $R^2 = 0.9952$
    ($RMSE = 0.2635$, $MAE = 0.1797$ in physical space), confirming that the learned correction drastically
    reduces the discrepancy with the LES-equivalent closure.

    \item The comparison with an equivalent network trained without
    clustering showed that the clustered model attains slightly better
    metrics across the board ($R^2 = 0.9952$ against $0.9950$, with lower
    $RMSE$ and $MAE$), reflecting a better reconstruction of the data
    variance. The benefit of clustering is
    most evident in the near-terrain region, where the clustered model
    is markedly more consistent with the LES reference, since it is able
    to learn separately the distinct near-wall and free-stream flow
    regimes.

    \item Once embedded in the solver, the ML-based RANS model
    consistently outperformed the standard RANS model in both the near-
    and far-wake. On the fine grid the model strongly improved the
    near-wake prediction, especially in the upper part of the wake,
    while slightly underestimating the deficit in the lower near-wake
    and overestimating it in the far-wake after the wake breakdown.

    \item The most relevant result for practical applications was
    obtained on the ultra-coarse grids: the model retained its improved
    near-wake prediction while also accurately reproducing the
    post-breakdown far-wake recovery, at a computational cost about
    $1.25$ times lower than the reference LES on the
    $\sim\!4.2\times10^6$ cells grid. Since downstream turbines are
    typically located in the far-wake of the upstream machines, this
    combination of far-wake accuracy and reduced cost is precisely what
    is required for wind farm layout and energy-yield studies.

    \item Although trained on a single inflow velocity
    ($U_\infty = 11\,\mathrm{m/s}$), the closure generalised well to the
    unseen $8$ and $15\,\mathrm{m/s}$ conditions: a scalar metric based on
    the integral of the velocity-deficit difference with respect to the
    LES showed that, at all three inflow velocities, the standard RANS
    model systematically overestimates the wake deficit while every
    RANS-ML configuration remains close to the LES reference throughout
    the wake.

    \item Beyond the velocity field, the model also reproduced the rotor
    forces and the $C_p$/$C_t$ coefficients in good agreement with the
    reference curves, with only a slight underestimation partly
    attributable to the more realistic ABL and ground-affected operating
    conditions, demonstrating its ability to capture the aerodynamic
    loading and the associated power production.
\end{itemize}

Two aspects deserve to be stated clearly, so that they are not mistaken
for weaknesses. First, training on a single operating condition and a
compact dataset is a deliberate design choice and the very premise of the
methodology, not a constraint: the aim is to calibrate the closure once,
on an affordable amount of high-fidelity data, and to reuse it across the
operating envelope, as the transfer to the unseen $8$ and
$15\,\mathrm{m/s}$ conditions and to coarser grids confirms. The mild
train/validation gap is the expected signature of this compact
calibration and, as shown, it does not degrade the out-of-sample
performance. Second, by construction the model remains a scalar
eddy-viscosity closure: it corrects the magnitude and distribution of
$\nu_t$ to reproduce the LES mixing, but does not represent the
Reynolds-stress anisotropy, a deliberate trade-off that keeps the closure
directly embeddable in standard solvers.

Genuine residual limitations remain and outline the path forward. A
residual underestimation bias persists in the far-upstream inflow region,
inherited from the synthetic inlet boundary condition, and the fine-grid
configuration still overestimates the far-wake deficit. Future work will
therefore enlarge the training dataset to multiple inflow velocities,
turbulence intensities and thermal-stability conditions, so as to broaden
the validated operating range and further tighten the closure; it will
explore tensorial, anisotropy-aware corrections to overcome the scalar
limitation; and it will extend the methodology to multi-turbine
configurations to assess wake-interaction effects. The compactness and
the reduced cost of the proposed data-driven closure also make it a
promising candidate for future real-time wind farm monitoring and control
applications, where accurate yet inexpensive representations of the wake
field are essential.
